% !PRINTED 121
\weberpage{121}
Wenn nun $e$ ein beliebiges Element aus $Q$ ist, so ist $Qe = Q$,
und folglich auch $Ra_iQe = Ra_iQ$. Daraus folgt, dass durch
Composition mit $e$ die Gesammtheit der $h_i$ Elemente, aus denen
$P_i$ zusammengesetzt ist, nicht geändert wird, und dass sie also
durch die Permutationsgruppe $\Pi$ nur unter einander vertauscht
werden. Man sieht aber auch unmittelbar, dass diese $h_i$ Elemente $B$ durch $\Pi$ \gesperrt{transitiv} verbunden sind; denn $P_1$ besteht eben
aus allen den Complexen $B$, die aus einem unter ihnen durch
Anwendung von Elementen aus $Q$ abgeleitet werden können, und
daher kann jeder dieser Complexe $B$ aus $P_i$ durch Permutationen
aus $\Pi$ in jeden anderen übergeführt werden.

Wir wollen das gewonnene Resultat in folgender Weise als
Theorem formuliren:

3.\ \gesperrt{Ist $P$ eine Gruppe vom Grade $n$, $R$ ein Theiler
von $P$ vom Index $j$, $Q$ ein Theiler von $P$ vom
Grade $q$, so ist mit $Q$ eine Permutationsgruppe $\Pi$
von $j$ Ziffern ein- oder mehrstufig isomorph. Die
$j$ Ziffern zerfallen in Reihen von je $h_1, h_2, h_3, \ldots$
durch $\Pi$ transitiv verbundenen Ziffern, so dass}
\weq{11}{j = h_1 + h_2 + h_3 + \cdots}
\gesperrt{ist, und $h_1, h_2, h_3, \ldots$ Theiler von $q$ sind. Die
Zahlen $h_1, h_2, h_3, \ldots$ sind die Indices von Theilern
von $Q$, die sich als Durchschnitte von $Q$ mit
den conjugirten Gruppen $a^{-1}Ra$ ergeben.}

Haben die verschiedenen Gruppen $Q_1, Q_2, \ldots$ ausser dem
Einheitselemente keinen gemeinschaftlichen Theiler, was bei einfachen Gruppen $P$ immer eintritt, so ist der Isomorphismus einstufig.

\begin{center}
\S.\,29.

\gesperrt{Der erste \gesperrt{Sylow}'sche Satz.}
\end{center}

Es ist nicht leicht, auf dem von \gesperrt{Cayley} eingeschlagenen
directen Wege bei der Bildung von Gruppen über die niedrigsten
Gradzahlen hinauszugehen. Für weitergehende Untersuchungen
in dieser Richtung ist ein Satz von grossem Nutzen, der in beschränktem Umfange von \gesperrt{Cauchy} herrührt, allgemein aber
von \gesperrt{Sylow} bewiesen ist. Wir wollen diesen wichtigen Satz
%CONT>>
